% Generated draft: requires mathematical and source review.
\documentclass[11pt]{article}
\usepackage[margin=25mm]{geometry}
\usepackage{fontspec}
\setmainfont{Latin Modern Roman}
\usepackage{amsmath,amssymb,mathtools,unicode-math}
\setmathfont{Latin Modern Math}
\usepackage{enumitem,needspace,xurl,hyperref}
\setlength{\parindent}{0pt}
\setlength{\parskip}{5pt}
\setlength{\emergencystretch}{3em}
\newcommand{\N}{\mathbb{N}}
\newcommand{\Z}{\mathbb{Z}}
\newcommand{\Q}{\mathbb{Q}}
\newcommand{\R}{\mathbb{R}}
\newcommand{\C}{\mathbb{C}}
\newcommand{\F}{\mathbb{F}}
\newcommand{\A}{\mathbb{A}}
\newcommand{\PP}{\mathbb{P}}
\newcommand{\E}{\mathbb{E}}
\newcommand{\eps}{\varepsilon}
\newcommand{\dd}{\,\mathrm{d}}
\begin{document}
\section*{UnsolvedMath ID: 11000311}
\textbf{Source problem code:} catalogue entry
\subsection{Definitions and mathematical statement}
In the standard convention let
\[
 I\operatorname{Out}_n=
 \ker\bigl(\operatorname{Out}(F_n)\to\operatorname{GL}(H_1(F_n;\Z))
 \cong\operatorname{GL}_n(\Z)\bigr).
\]
The source does not provide this definition or the construction of its higher Franz--Reidemeister torsion classes
\[
 \tau_{2i}\in H^{4i}(I\operatorname{Out}_n;\R).
\]
Here ``torsion'' names a secondary characteristic invariant; it does not mean a finite-order element of real cohomology.

\emph{Problem 4.4 (Igusa).} Prove that these classes are nonzero in a suitable stable range. In the usual positive-index reading, for every fixed \(i\geq1\), determine a suitable range of ranks \(n\) in which \(\tau_{2i}\neq0\); an eventual-rank formulation would assert
\[
 \forall i\geq1\ \exists N(i)\ \forall n\geq N(i),
 \qquad \tau_{2i}(n)\neq0.
\]
The actual stable range, construction and normalization, and rank dependence of the classes are not supplied. The displayed quantifiers explain the conventional stable-range goal without inventing a numerical threshold or rank-stabilization maps.

The source then discusses an unstable regime in which relevant Borel classes vanish in \(H^*(\operatorname{GL}_n(\Z);\R)\). It suggests relations between \(\tau_{2i}\), the duals of classes denoted \(\mu_i\), and unstable rational cohomology of \(\operatorname{GL}_n(\Z)\). The \(\mu_i\), degrees, vanishing range, and relations are not defined; the wording is conjectural. The following specific example promised by the source is absent and is not reconstructed.
\subsection{Short English statement}
Igusa asks for nonvanishing of the real higher Franz--Reidemeister torsion classes \(\tau_{2i}\) in degree \(4i\) for a suitable stable rank range. The neighboring unstable Borel/\(\mu_i\) discussion is conjectural context, not a supplied relation or proof.
\subsection{Sources}
\begin{itemize}
\item Local UnsolvedMath registry, permanent ID 11000311, source code catalogue entry, Problem 4.4, PDF page 365; explicit attribution Igusa.
\item Farb, \emph{Problems on Mapping Class Groups and Related Topics}. Supplied URL: \url{https://www.math.uchicago.edu/~farb/papers/mcgbook.pdf}; recorded access 2026-07-29; extraction pdf-font-structure. The volume contains 329 explicitly typeset headings across chapters, with repeated labels disambiguated by page. The chapter header is ``Cohomological structure of the mapping class group and beyond.''
\item No full citation for the torsion-class construction, \(\mu_i\), or Borel comparison is supplied. The summary reports partial/low-case results and says higher nonvanishing remains unproved to its author's knowledge, without precise references.
\item Supplied status open and OPEN-TRIAGE. Current primary-literature status and redistribution terms for the publicly accessible author/publisher copy remain NEEDS\_REVIEW.
\end{itemize}
\end{document}
